% Reference: https://www.luogu.com.cn/article/2hb00cd4 (non-polynomial portions).
% Table cross-check: https://dlmf.nist.gov/26.17
% Sources and validation: docs/validation.md.
\section{十二重计数：球与盒子}
\label{math:twelvefold}
将 $n$ 个球全部放入 $m$ 个盒子，不计同一盒内的顺序。球、盒子分别是否可区分，再结合“不限制、每盒至多一个、每盒至少一个”三类条件，得到十二种模型。先用下表定位，再按后面的计数方法计算；表中取 $n,m\geq1$。

\begin{center}
\small
\renewcommand{\arraystretch}{1.5}
\begin{tabularx}{\linewidth}{@{}ccXXX@{}}
\toprule
球 & 盒子 & 不限制 & 至多一个 & 至少一个 \\
\midrule
不同 & 不同 & $m^n$ & $\dfrac{m!}{(m-n)!}$ & $m!S(n,m)$ \\
相同 & 不同 & $\binom{n+m-1}{m-1}$ & $\binom mn$ & $\binom{n-1}{m-1}$ \\
不同 & 相同 & $\sum_{j=0}^m S(n,j)$ & $[n\leq m]$ & $S(n,m)$ \\
相同 & 相同 & $P_{\leq m}(n)$ & $[n\leq m]$ & $P_{\leq m}(n-m)$ \\
\bottomrule
\end{tabularx}
\end{center}
其中 $S(n,j)$ 是第二类 Stirling 数，$P_{\leq m}(n)$ 是至多分成 $m$ 部分的整数划分数；$[Q]$ 表示条件成立为 $1$，否则为 $0$。不可行的排列数、组合数和负参数的划分数统一记为 $0$。

\subsection*{球不同、盒子不同：分配与容斥}
每个球独立选择盒子，不限制时为 $m^n$；至多一个时，将 $n$ 个球依次分配到互不相同的盒子，为下降积 $m(m-1)\cdots(m-n+1)$，$n>m$ 时无解。

要求每盒非空时，以“某些盒子为空”为坏事件。选定 $j$ 个空盒后，所有球只能放入其余 $m-j$ 个盒子，故由容斥得到
\[
    F(n,m)=\sum_{j=0}^m(-1)^j\binom mj(m-j)^n=m!S(n,m).
\]
预处理组合数后，逐项快速幂可在 $O(m\log(n+1))$ 时间计算。也可按非空盒子数分类写成
\[
    m^n=\sum_{j=0}^m\binom mj F(n,j),
\]
它与上面的交错求和互为二项式反演：先选哪些盒子非空，再往其中满射。

\subsection*{球不同、盒子相同：集合划分}
$S(n,j)$ 计数将 $n$ 个不同元素分成 $j$ 个非空、无标号集合块。把最后一个元素单独成块，或放入已有的某一块，得到
\[
    S(n,j)=S(n-1,j-1)+jS(n-1,j).
\]
初值 $S(0,0)=1$，$S(n,0)=0$（$n>0$），$j>n$ 时为 $0$。递推至 $n$ 行、$\min(n,m)$ 列需要 $O(n\min(n,m))$ 时间，可滚动到 $O(\min(n,m))$ 空间；一维实现应倒序枚举 $j$，并在每一行更新后将第 $0$ 项置为 $0$。

允许空盒时，只关心非空集合块数，求 $\sum_{j=0}^m S(n,j)$；要求盒子非空时取 $S(n,m)$。至多一个时，每个球单独成块，若 $n\leq m$ 只有一种方案。

满射公式除去 $m$ 个非空集合块的标号，还得到
\[
    S(n,m)=\frac1{m!}\sum_{j=0}^m(-1)^j\binom mj(m-j)^n.
\]
在整数意义下该商为整数；取模时只有 $m!$ 可逆才能直接乘逆元，否则用不涉及除法的 Stirling 递推。一般情况下，不能把有标号盒子的答案简单除以 $m!$：空盒或相同内容可能造成不同数量的重复。

\subsection*{球相同、盒子不同：插板法}
令 $x_i$ 表示第 $i$ 个盒子的球数。若每盒至少一个，方程 $x_1+\cdots+x_m=n$、$x_i\geq1$ 的解可由球之间的 $n-1$ 个间隙选出 $m-1$ 块隔板，得到 $\binom{n-1}{m-1}$。

若允许为空，令 $y_i=x_i+1$，改为把 $n+m$ 个球正分成 $m$ 份，得到 $\binom{n+m-1}{m-1}$。若至多一个，只需选择哪 $n$ 个盒子放球，得到 $\binom mn$。

\subsection*{球相同、盒子相同：整数划分}
将各盒数量按非增序排列，相同的数量序列只计一次，所以不限制时答案为 $P_{\leq m}(n)$。按“有空盒”或“所有 $m$ 个盒子均非空”分类，后一类从每盒取走一个球，可得递推
\[
    P_{\leq m}(n)=P_{\leq m-1}(n)+P_{\leq m}(n-m).
\]
边界为 $P_{\leq m}(0)=1$、$P_{\leq0}(n)=[n=0]$，以及 $n<0$ 时为 $0$。要求非空时直接取 $P_{\leq m}(n-m)$；至多一个仍为 $[n\leq m]$。

把划分画成方格图并交换行列，可知“至多 $m$ 部分”等价于“每部分大小不超过 $m$”。因此也可做完全背包：初始 $dp[0]=1$，其他为零；先枚举部分大小 $i=1,\ldots,\min(n,m)$，再从 $s=i$ 到 $n$ 正序更新 $dp[s]\gets dp[s]+dp[s-i]$。最终 $dp[n]=P_{\leq m}(n)$，时间 $O(n\min(n,m))$，空间 $O(n)$。

\subsection*{边界与计算方式}
$n=0$ 时，允许空盒或每盒至多一个均只有空分配一种；要求每盒至少一个时，仅 $m=0$ 可行。$n>0,m=0$ 时无解。容斥式在空对象边界按 $0^0=1$ 解释，与唯一的空映射一致。

公式本身不限制模数。实际计算时，快速幂、阶乘表、Stirling 递推和划分 DP 按数据范围选择；取模的除法仍须先检查逆元是否存在。
